\chapter{序、序数与超限构造}
\begin{qlnote}{阅读参照}Enderton, \emph{Set Theory}, Chapters 7--8。
本章先固定定义与后续证明任务。\end{qlnote}

\section{偏序、线序与良序}
非严格偏序 \(\le\) 满足自反、反对称和传递；线序还要求任意两元素可比较。
严格序 \(<\) 的公理使用反自反性；不能把两套约定混在同一个定义里。
一个严格线序是\textbf{良序}，若每个非空子集都存在最小元素。
\(\omega\) 的通常序是良序；\(\mathbb Z\) 的通常序不是，因为负整数集合没有最小元素。

\section{Von Neumann ordinals}
\begin{definition}{Transitive sets and ordinals}
集合 \(A\) 传递，若 \(x\in y\in A\Rightarrow x\in A\)。
一个传递集合 \(\alpha\)，若其成员被 \(\in\) 严格良序，则称为\textbf{序数}。
对序数写 \(\alpha<\beta\) 当且仅当 \(\alpha\in\beta\)。
\end{definition}
有限自然数和 \(\omega\) 是标准例子。序数后继为 \(\alpha+1=\alpha\cup\{\alpha\}\)；
非零且不是后继的序数称为极限序数。
成员关系与子集关系在序数之间有特殊联系，但不能把这种联系推广到任意集合。

\section{超限归纳与递归的角色}
对一个固定序数 \(\alpha\)，若 \(A\subseteq\alpha\) 满足：对每个 \(\beta<\alpha\)，
只要所有 \(\gamma<\beta\) 属于 \(A\)，就有 \(\beta\in A\)，则 \(A=\alpha\)。
否则 \(\alpha\setminus A\) 的最小元素给出矛盾。

超限递归在阶段 \(\beta\) 根据先前整段函数 \(f\restriction\beta\) 定义新值。
它需要精确说明可定义的规则、值的唯一性及集合长度的定义域。
替代公理用于把集合索引的已定义值收集成集合；极限阶段还需证明近似图像兼容。
序数可比较性、每个良序的序型、超限递归及 rank 的完整构造，随原书逐项展开。

\section{序型与大小}
\(\omega+1\) 的序是在自然数之后添加一个最大元素，所以它不与 \(\omega\) 序同构，
尽管两者的底层集合等势。序数描述良序的型；基数描述集合的大小。
在选择公理下，每个集合可良序，其大小可由某个最小的等势序数，即初始序数表示。
集合的全体与所有序数的全体要区别：后者是一个真类，不能未经论证放进幂集或作为函数的集合定义域。
